\begin{abstract}
Cybercrime in India has grown rapidly across diverse linguistic and regional contexts,
yet existing detection systems remain limited to monolingual, text-based analysis.
This paper presents a transformer-based multilingual multimodal cybercrime detection
framework that jointly analyzes textual and visual inputs — including complaint text,
screenshots, and fraudulent interfaces — collected from official Indian cybercrime
portals. We introduce a dataset of 25,000 annotated samples across English, Hindi,
Telugu, Tamil, and Malayalam, organized into a hierarchical taxonomy of 50 cybercrime
categories spanning financial fraud, identity theft, social engineering, and emerging
AI-based threats. The proposed framework encodes text using Qwen2.5-7B and visual
content using ViT-B/16, fusing both modalities via cross-modal attention. Explainability
is provided through LIME for textual inputs and attention maps for visual content.
The model achieves a macro F1 of 0.883 and accuracy of 0.892, outperforming the
strongest text-only baseline (XLM-R) by 5.3 F1 points and demonstrating consistent
performance across all five languages (macro F1: 0.869). The framework offers a
scalable, interpretable solution for automated cybercrime analysis in the Indian context.
\end{abstract}